% Generated draft: requires mathematical and source review.
\documentclass[11pt]{article}
\usepackage[margin=25mm]{geometry}
\usepackage{fontspec}
\setmainfont{Latin Modern Roman}
\usepackage{amsmath,amssymb,mathtools,unicode-math}
\setmathfont{Latin Modern Math}
\usepackage{enumitem,needspace,xurl,hyperref}
\setlength{\parindent}{0pt}
\setlength{\parskip}{5pt}
\setlength{\emergencystretch}{3em}
\newcommand{\N}{\mathbb{N}}
\newcommand{\Z}{\mathbb{Z}}
\newcommand{\Q}{\mathbb{Q}}
\newcommand{\R}{\mathbb{R}}
\newcommand{\C}{\mathbb{C}}
\newcommand{\F}{\mathbb{F}}
\newcommand{\A}{\mathbb{A}}
\newcommand{\PP}{\mathbb{P}}
\newcommand{\E}{\mathbb{E}}
\newcommand{\eps}{\varepsilon}
\newcommand{\dd}{\,\mathrm{d}}
\begin{document}
\section*{UnsolvedMath ID: 20001924}
\textbf{Source problem code:} catalogue entry
\subsection{Definitions and mathematical statement}
Question 9.1.1 asks whether every closed Lorentzian manifold has at least one closed geodesic. The supplied verification specifies the source convention: a nonconstant simple closed geodesic is an embedding \(\gamma:S^1\to M\) whose periodic affine parametrization satisfies
\[
 \nabla_{\dot\gamma}\dot\gamma=0.
\]
Lorentzian means a nondegenerate metric of index one, up to the opposite sign convention. Closed means compact without boundary and does not imply geodesic completeness in this signature. A weakly closed geodesic image, whose affine velocity returns only up to a nonunit positive factor, is different from an affine periodic orbit.

The aid reports a 2024 negative answer by Allout, Belkacem and Zeghib: compact three-dimensional Lorentzian quotients associated with \(SL(2,\mathbb R)\) or \(\mathrm{Sol}\) have no affine closed geodesic. Weakly closed incomplete null geodesics remain in those examples. The record supplies no quotient construction or full bibliographic entry. The reported counterexamples therefore must not be rephrased as excluding all weakly closed images or as settling every geodesically complete variant.

The source also asks about arbitrary signatures and finite-volume complete manifolds, expressly warning that Lorentzian completeness has several inequivalent meanings. It cites the introduction of [Fl-Ja-Pi2011] for known closed timelike-geodesic results under additional assumptions, without listing those assumptions. The tail begins Section 9.2, ``Light-like geodesics of semi-Riemannian metrics,'' on a closed manifold with indefinite metric; it contains no complete additional question here.

The paired attempt gives a separate positive symmetry criterion for a compact semi-Riemannian manifold of any nondegenerate signature. A nontrivial isometric \(S^1\)-action forces an affine simple closed geodesic. If \(X\) is its Killing generator and \(f=g(X,X)\), the underlying critical-point criterion uses
\[
 \nabla_XX=-\tfrac12\operatorname{grad}f.
\]
It does not assert that the forced geodesic is timelike.

For a free isometric circle action, \(B=M/S^1\) is a smooth compact orbit space and \(f\) descends to \(\overline f:B\to\mathbb R\). The aid reports distinct simple affine geodesic circle orbits over critical points of \(\overline f\). If \(\overline f\) is Morse, then for any coefficient field \(\mathbb F\),
\[
 \#\{\text{geometrically distinct geodesic }S^1\text{-orbits}\}
 \geq\#\operatorname{Crit}(\overline f)
 \geq\sum_j b_j(B;\mathbb F).
\]
Freeness and the Morse assumption are essential to this stated refinement. A final reported consequence is that a compact semi-Riemannian manifold with no affine closed geodesics has only finite compact subgroups of its isometry group: a positive-dimensional compact Lie subgroup would contain a nontrivial circle subgroup. This is a restriction on compact subgroups, not a claim that the whole isometry group is finite.
\subsection{Short English statement}
The supplied report attributes a 2024 Lorentzian counterexample to Allout, Belkacem and Zeghib, with affine and weakly closed geodesics carefully distinguished. A separate criterion forces affine closed geodesics from an isometric circle action, with a Morse lower bound for free actions. The source's completeness variants and adjacent lightlike heading remain separately qualified.
\subsection{Sources}
\begin{itemize}
\item Local source record: catalogue entry, catalogue entry; canonical \texttt{aim-geometry-notes.json}, \texttt{notes[261]}, source item 9.1, Question 9.1.1. Title: ``The 2024 Lorentzian refutation and a proved symmetry-forced Morse criterion.''
\item Keith Burns and Vladimir S. Matveev's AIM list \emph{Open Problems and Questions about Geodesics}, workshop ``Geodesics,'' supplied reference \url{https://aimath.org/pastworkshops/geodesicsproblems.pdf}. This draft uses the supplied record without accessing a PDF. Its verification supplies the affine simple-closed convention; the page header is an extraction artifact.
\item Original reference [Fl-Ja-Pi2011], introduction, for timelike results under additional assumptions. Supplied report attribution: Allout, Belkacem and Zeghib, 2024, for the compact Lorentzian refutation. No full bibliographic entries are supplied.
\item Attempt 1 supplies the circle-action and Morse-count criteria. Registry dated 14 August 2026: ``solved,'' ``SOLVED-IN-LITERATURE,'' statement status ``exact''; research-run status refuted, attempt partial, rigor ``proved,'' novelty confidence low. The machine-generated, non-peer-reviewed aid requires independent verification before citation.
\end{itemize}
\end{document}
